\section{最上次数テンソルによるFrobenius候補の復元}
\label{sec:frobenius-tensor}
この節ではKoszul性を仮定する前の有限Frobenius候補を整理する。
最上次数のテンソルを用い、正則性の未決定と、全対象を覆う不変な分岐を分ける。

\begin{proposition}[Frobenius検査の四次への短縮]\label{prop:frobenius-degree-four}
四窓の二次双対$E$について、各始点の次数$0,\ldots,4$の次元が
$(1,4,6,4,1)$であり、相補次数の積対合が完全なら、$E_{i,i+5}=0$は自動的に成立する。
従って命題\ref{prop:finite-frobenius}の五次元数検査はこの仮定の下で省略できる。
これは元の代数のKoszul性の判定を短縮する主張ではない。
\end{proposition}
\begin{proof}
前節の交差空間$W_{i,p}=E_{i,i+p}^*$を用いる。
非零traceを引き戻したテンソル$w_i\in T_{i,4}$は$W_{i,4}$を生成する。
$E_{i,i+1}\otimes E_{i+1,i+4}$の完全対合から、
$w_i$の最初の一因子と残りの三因子の間の行列階数は4である。
一方、交差空間の定義により
\[
 W_{i,5}=(W_{i,4}\otimes V_{i+4})\cap(V_i\otimes W_{i+1,4}).
\]
非零元があれば、ある非零$u\in V_i$, $v\in V_{i+4}$により
$w_i\otimes v=u\otimes w_{i+1}$と書ける。
$v$を1へ送る線型汎関数で末尾を縮約すると$w_i=u\otimes z$となり、
最初の一因子での行列階数が1となって矛盾する。
よって$W_{i,5}=0$であり、次数一生成性から全高次も零となる。
\end{proof}

$0\ne w\in V_0\otimes V_1\otimes V_2\otimes V_3$を取り、
その巡回回転を$w_i\in V_i\otimes\cdots\otimes V_{i+3}$と書く。
回転には符号を入れず、添字は法4で読む。
テンソルを二組の因子に分けた線型写像の階数を、その分け方での行列階数と呼ぶ。
$w_i$の後ろ二因子を縮約して得る像を
\[
 R_i(w)=\operatorname{im}\bigl((V_{i+2}\otimes V_{i+3})^*
                  \xrightarrow{\ w_i\ } V_i\otimes V_{i+1}\bigr)
\]
と定め、これらから交差空間$W_{i,p}(w)$を作る。

\begin{theorem}[標準Frobenius四窓のテンソル復元]\label{thm:frobenius-tensor-reconstruction}
射影テンソル$[w]$に次の有限条件を課す。
\begin{enumerate}
\item 各一因子と残り三因子の間の行列階数は4である。
\item $01\mid23$と$12\mid30$の二つの行列階数は6である。
\item 全ての始点で$\dim W_{i,3}(w)=4$, $\dim W_{i,4}(w)=1$である。
\end{enumerate}
これらの$[w]$と、有限Frobenius検査を通り厳密4移動がFrobenius標準である四窓とは、
$[w]\mapsto(R_i(w))$によって$G$同変な全単射となる。
従って、さらに元代数のKoszul AS正則性を課した部分では、
$G$軌道が裸の標準14641型の添字保持同型類を与える。
上記三条件だけでKoszul性を結論することはしない。
\end{theorem}
\begin{proof}
まず三条件を満たす$w$から出発する。
二つの二対二階数条件により各$R_i(w)$は6次元である。
テンソル$w_i$は、隣接する三つの位置の関係条件全てを満たすので$W_{i,4}(w)$に入り、
第三条件によりその生成元になる。
$w_i$の最初または最後の一因子を縮約した像は、対応する三次交差空間に含まれる。
第一条件と第三条件により、これは4次元の三次交差空間全体に一致する。
$E$をこの関係四窓の二次双対とし、$w_i$を引き戻しとして持つtrace
$\lambda_i:E_{i,i+4}\to k$を選ぶ。
一対三の積対合は第一条件から、二対二の積対合は第二条件から完全であり、
零対四と四対零も完全である。
前命題から五次以上は零なので、有限Frobenius条件が全て成立する。
回転の定義そのものにより
\[
 \lambda_i(\alpha\beta\gamma\delta)
   =\lambda_{i+1}(\beta\gamma\delta\alpha)
\]
が成立する。ここで各文字は対応する次数一双対空間の元である。
次数三の積が$E_{i+1,i+4}$を生成するから、この等式は局所Nakayama写像が
厳密4移動であることを意味する。

逆に標準Frobenius四窓から出発し、標準性を与えるtraceを選ぶ。
その引き戻し$w_i$は$W_{i,4}$を生成し、Nakayamaの特徴付けから互いに巡回回転である。
一対三の完全対合により第一条件が成立する。
二対二の完全対合から、後ろ二因子の縮約像は$E_{i,i+2}^*=R_i$全体になり、
第二条件と関係空間の復元を得る。第三条件は双対の次数元数である。
$W_{0,4}$が一次元なので$[w]$は一意である。
全構成は基底変更と可換し、最後の同型類の主張は系\ref{cor:canonical-slice}による。
\end{proof}

\begin{proposition}[第三の二対二階数による全被覆分岐]\label{prop:opposite-rank-split}
標準14641型の対象について、上の最上次数テンソルから
\[
 \rho(C)=\rk\bigl((V_0\otimes V_2)^*\xrightarrow{\ w\ }V_1\otimes V_3\bigr)
\]
を定めると、これはtrace・基底・標準四窓の選択によらない添字保持同型不変量であり、
$2\leq\rho(C)\leq16$を満たす。
従って$\rho\leq6$と$\rho\geq7$は、全対象を互いに素な二つの部分へ分ける。
前者はテンソル候補空間内で相対的に閉、後者は相対的に開である。
各部分の正規形・AS正則locus・完全同型判定はなお未決定である。
\end{proposition}
\begin{proof}
traceの変更は$w$の非零スカラー倍、基底変更は四因子での可逆変換を与えるため階数を保つ。
裸の同型を標準構造を保つ同型へ補正できることは定理\ref{thm:canonical-four}で示した。
よって標準四窓の選択にもよらない。
非零$16\times16$行列なので$1\leq\rho\leq16$である。
もし$\rho=1$なら$w=a_{02}\otimes b_{13}$と書ける。
四つの一因子階数が4だから、$a_{02},b_{13}$を$4\times4$行列と見た階数はともに4。
すると$01\mid23$での階数はそのKronecker積の階数16となり、6という条件に矛盾する。
二分岐の被覆と非交差性は整数階数の二者択一であり、
閉性は7次小行列式の消滅による。
\end{proof}
\begin{remark}
これは正則な全対象に対する被覆分岐であって、各子の分類が完成したという主張ではない。
同じ$\rho$でも同型とは限らず、異なる$\rho$なら非同型である。
定理のテンソル候補族の中の特殊化では階数は増加しない。
\path{computations/verify_frobenius_tensor.py}は復元と既存族の階数を有理数上で検査する。
テンソル表示の一般原理そのものの新規性は主張しない。
\end{remark}

\subsection{量子多項式上の全階数層と最小階数層の同型}
\begin{proposition}[量子族の第三階数]\label{prop:quantum-opposite-rank}
量子多項式代数
\[
 A(q)=k\langle x_0,x_1,x_2,x_3\rangle/
 (x_bx_a-q_{ab}x_ax_b:\ 0\leq a<b\leq3),\qquad q_{ab}\in k^\times
\]
に対して
\[
 F_1=q_{02}q_{13}-q_{03}q_{12},\quad
 F_2=q_{01}q_{12}q_{23}-q_{03},\quad
 F_3=q_{01}q_{13}-q_{02}q_{23}
\]
と置く。このとき
\[
 \rho(\check A(q))=6+2\#\{j\in\{1,2,3\}:F_j\ne0\}.
\]
従って量子族で現れる階数は正確に$6,8,10,12$であり、一般点では12である。
$\rho=6$の量子対象は、添字保持同型でちょうど二種類、
可換多項式代数と全ての$q_{ab}=-1$である反可換量子多項式代数に帰着する。
これは量子族に限定した主張であり、全対象の$\rho=6$層の分類ではない。
\end{proposition}
\begin{proof}
量子書換えは全重なりを解消し、順序付き単項式が基底になる。
二次双対は量子外積代数で完全対合を持つため、全ての$q_{ab}\ne0$で
標準型Koszul AS正則である。
四次交差空間の生成元は、$\pi\in S_4$に対し
\[
 w=\sum_{\pi\in S_4}c_\pi\,
 x_{\pi(0)}\otimes x_{\pi(1)}\otimes x_{\pi(2)}\otimes x_{\pi(3)},\qquad
 c_\pi=\prod_{a<b}
 \begin{cases}q_{ab},&\pi^{-1}(a)<\pi^{-1}(b),\\-1,&\pi^{-1}(a)>\pi^{-1}(b)\end{cases}
\]
で与えられる。隣接する$a,b$を交換した二係数の比が$-q_{ab}$になるので、
各隣接位置の関係条件を満たすことが直接分かる。
このテンソルは標準化前の最初の四一次空間だけで定まる。
それらの基底を保ったまま標準4周期を延長しても$w$は変わらないので、$\rho$の計算に使える。

$02\mid13$行列の非零行は異なる二文字を持つ12行である。
行の二文字集合$S\subset\{0,1,2,3\}$を固定すると、列の二文字集合はその補集合に限られる。
従って六つの$2\times2$ブロックに分解する。
全ブロックの各成分は非零単項式であり、行集合$S$を順に
$01,02,03,12,13,23$と取った行列式は、可逆単項式を除いてそれぞれ
$F_1,F_2,F_3,F_3,F_2,F_1$である。
よって各ブロックの階数は行列式零で1、非零で2となり公式を得る。
これらの恒等式は\path{computations/verify_quantum_tensor_rank.py}でも記号的に検査する。
例えば$q_{01}=q_{12}=q_{23}=1$と固定し、
$(q_{02},q_{13},q_{03})=(1,1,1),(2,2,1),(2,3,1),(2,3,5)$と取れば、
それぞれ階数$6,8,10,12$を得る。非零条件は開であり一般点の主張も従う。

階数6なら$F_1=F_2=F_3=0$なので
\[
 q_{03}=q_{01}q_{12}q_{23},\qquad
 q_{13}=q_{02}q_{23}/q_{01},\qquad
 q_{02}^2=q_{01}^2q_{12}^2.
\]
標数零より、ある$s\in\{1,-1\}$に対して
$q_{02}=s q_{01}q_{12}$, $q_{13}=s q_{12}q_{23}$である。
$d_0=1$, $d_1=q_{01}/s$, $d_2=q_{01}q_{12}$,
$d_3=q_{01}q_{12}q_{23}/s$と置くと
$q'_{ab}=q_{ab}d_a/d_b=s$が全ての対で成立する。
位置$i$の基底変換$g_i(x_a)=d_a^i x_a$は、
$q'$の関係を$q$の関係へ送る。従ってこれは全整数の添字を保つ同型であり、
graded同型を主張するものではない。

最後に二つの標準形が添字保持同型でないことを示す。
最初の二次関係の左点対応の射影支持は、可換型では全$\mathbb P^3$である。
反可換型では、$x_ay_b+x_by_a=0$を満たす非零$y$が存在するためには、
$x$の非零座標は高々二つでなければならない。
実際、$x_a\ne0$なら$y_a=0$は全$y=0$を強制し、三つの非零座標があれば
三対の式が$y_b/x_b=-y_a/x_a$を互いに矛盾させる。
逆に非零座標が高々二つならそのような$y$を選べる。
従って支持は六本の座標直線の和である。
添字保持同型は二一次空間の射影変換によりこの支持を保つので、両型は非同型である。
\end{proof}
